\documentclass[10pt,a4paper]{amsart}
\usepackage[margin=19mm]{geometry}
\usepackage{amsmath,amssymb,mathtools}
\usepackage[scr=rsfs,cal=euler]{mathalfa}
\usepackage{xfrac}
\usepackage[unicode,hidelinks,bookmarks=false]{hyperref}
\newcommand{\ie}{i.e.\ }
\newcommand{\cf}{cf.\ }
\newcommand{\resp}{respectively\ }
\newcommand{\Subsection}{\subsection}
\providecommand{\nobreakdash}{\nobreak\hbox{-}}
\setlength{\emergencystretch}{2em}
\multlinegap0pt
\usepackage{bm}
\usepackage{amssymb}
\usepackage{xcolor}
\usepackage[normalem]{ulem}
\usepackage{algorithm}\usepackage{algpseudocode}
\newtheorem{assumption}{Assumption}
\newtheorem{theorem}{Theorem}
\newcommand{\R}{\mathbb{R}}
\newcommand{\dd}{\mathrm{d}}
\newcommand{\op}{\operatorname}
\title{Dual Approaches to Stochastic Control via SPDEs and the Pathwise Hopf Formula}
\author{Mathieu Lauri\`ere, Jiefei Yang}
\date{Source: arXiv:2604.08155v1 (CC BY 4.0)}
\begin{document}
\maketitle

\noindent\emph{Source and licence.} Dual Approaches to Stochastic Control via SPDEs and the Pathwise Hopf Formula. Version arXiv:2604.08155v1 (CC BY 4.0), distributed by arXiv under the Creative Commons Attribution 4.0 International licence, http://creativecommons.org/licenses/by/4.0/, as recorded in the arXiv licence metadata for this exact version and shown on the abstract page https://arxiv.org/abs/2604.08155v1. Full licence notice: \texttt{licenses/arxiv-2604.08155-CC-BY-4.0.txt}.\par\smallskip

\noindent\textit{Editorial scope.} Counted unit: the article's theorem thm:generalized-hopf (source0.tex lines 492--502). The article's complete original proof of it is delivered with it (source0.tex lines 503--658, one \texttt{proof} environment of 156 lines). No mathematics is written, repaired or re-derived: the statement and the proof are the article's own lines, in the article's own notation, and no retained line is substituted or re-flowed.  Retained uncounted as the source-local context the proof needs: lines 149--156; lines 348--353; lines 378--395; lines 397--418; lines 407--413; lines 468--490.  Nothing was omitted from any retained span, so the package carries no omission record.  No retained line contains a citation, so the wrapper carries no bibliography and no bibliography entry.  No retained line contains a figure environment, an image-inclusion command or drawing code, so no image is delivered and the package declares no asset dependency.  Editorial formatting only, in addition: an AMS \texttt{amsart} preamble that restates the article's own declarations and macros used by the retained lines, namely the environments \texttt{assumption}, \texttt{theorem} on separate counters, together with the standard packages those lines need.  The retained environments print in this document as Assumption 1, Theorem 1 in order of appearance, whereas the article prints the same items with the numbering of its own full document.  The complete native source of the article as distributed in the arXiv e-print for this exact version is bundled unchanged as \texttt{sources/198139/source0.tex}.
\begin{assumption} \label{assu:dual-sc-prob}
    Assume that 
    \begin{enumerate}
        \item $b, r$ are uniformly bounded and continuous in $u$. 
        \item $b, \sigma, r$ are uniformly Lipschitz in $x$. 
        \item $g$ is Lipschitz continuous. 
    \end{enumerate}
\end{assumption}

\begin{equation} \label{eq:doc-prob}
\begin{aligned}
    v_n(0,x_0) =&\inf_u \bigg\{ \int_0^T \bigg(r(x(t), u(t)) - \langle z_\theta(t, x(t)), \dot{w}_n(t)\rangle + \frac{1}{2}\op{Tr}\left( \sigma(x) J_x z_\theta(t, x(t)) \right)\bigg)\,\dd t + g(x(T))\bigg\}, \\
    &\text{s.t. } \frac{\dd x(t)}{\dd t} = b(x(t), u(t)) - \frac{1}{2}(\nabla \sigma:\sigma)(x) + \sigma(x) \dot{w}_n(t). 
\end{aligned}
\end{equation}

\begin{assumption} \label{assu:Pontryagin}
Let $\bm{\xi}\in \R^{nd}$. Assume that 
\begin{enumerate}
    \item For any parameter $\theta$, the functions $(t,x)\mapsto \langle z_\theta(t,x), \dot{w}_n^{\bm{\xi}}(t)\rangle$ and $(t,x)\mapsto \op{Tr}(\sigma(x)J_xz_\theta(t,x))$ are uniformly continuous with respect to $(t,x)\in [0,T]\times \R^d$. The function $(\nabla \sigma:\sigma)(\cdot)$ is uniformly continuous with respect to $x\in \R^d$.
    \item There exists an $x_0\in \R^d$ and positive constant $A_1, A_2$ such that 
    \begin{equation*}
        \|b(x_0,u)\|_{\R^d} \le A_1, \|r(x_0,u)\|\le A_2  \quad \forall u\in U.
    \end{equation*}
    \item The functions $b(\cdot, u)$, $r(\cdot,u)$, $\sigma$, $g$, $z_\theta(t,\cdot)$, $J_xz_\theta(t,\cdot)$, and $(\nabla \sigma:\sigma)$ are continuously differentiable with respect to $x\in \R^d$ for any $t\in [0,T]$. 
    \item The function $\nabla g$ is uniformly continuous in $x$. $\nabla_x \tilde{b}(t,x,u; w_n^{\bm{\xi}})$ and $\nabla_x \tilde{r}(t,x,u; w_n^{\bm{\xi}})$ are uniformly continuous in $x$ and $u$ for any $t\in [0,T]$, where 
    \begin{equation} \label{eq:defi-tilde-b-r}
    \begin{aligned}
        \tilde{b}(t,x,u; w_n^{\bm{\xi}}) &:= b(x, u) - \frac{1}{2}(\nabla \sigma:\sigma)(x) + \sigma(x) \dot{w}_n^{\bm{\xi}}(t), \\
        \tilde{r}(t,x,u; w_n^{\bm{\xi}}) &:= r(x, u) - \langle z_\theta(t, x), \dot{w}_n^{\bm{\xi}}(t)\rangle + \frac{1}{2}\op{Tr}\left( \sigma(x) J_x z_\theta(t, x) \right). 
    \end{aligned}
    \end{equation}
\end{enumerate}
\end{assumption}

\begin{theorem}[Pontryagin's maximum principle]
\label{thm:pontryagin}
Let $\bm{\xi}\in \R^{nd}$. Let Assumptions \ref{assu:dual-sc-prob} and \ref{assu:Pontryagin} hold. Let $u^*(\cdot)$ be an optimal control for the problem~\eqref{eq:doc-prob} with $w_n = w_n^{\bm{\xi}}$ and $x^*(\cdot)$ be the optimally controlled state trajectory. Then there exists a function $p^*:[0,T] \to \R^d$ such that $(x^*, p^*)$ is a solution pair to the two point boundary value problem 
\begin{equation} \label{eq:two-point-boundary-value-prob}
    \begin{aligned}
        \dot{x}^*(t) &= \partial_p H(t, x^*(t), u^*(t), p^*(t); w_n^{\bm{\xi}}), \quad x^*(0) = x_0, \\
        \dot{p}^*(t) &= -\partial_x H(t,x^*(t), u^*(t), p^*(t); w_n^{\bm{\xi}}), \quad p^*(T) = \nabla g(x^*(T)), 
    \end{aligned}
\end{equation}
where we denote 
\begin{equation} \label{eq:hamiltonian-defi}
\begin{aligned}
    H(t,x,u,p; w_n^{\bm{\xi}}) := &\langle p, b(x,u)\rangle + r(x, u) + \left( p^\top \sigma(x) - z_\theta(t,x) \right) \dot{w}_n^{\bm{\xi}}(t)  \\
    &+ \frac{1}{2}\op{Tr}(\sigma(x) J_xz_\theta(t,x)) - \langle p, \frac{1}{2}(\nabla \sigma:\sigma)(x)\rangle, \\
    \mathcal{H}(t,x,p;w_n^{\bm{\xi}}) :=&\inf_{u\in U} H(t,x,u,p; w_n^{\bm{\xi}}), 
\end{aligned}
\end{equation}
and $u^*$ satisfies 
\begin{equation} \label{eq:u-star}
    H(t, x^*(t), u^*(t), p^*(t); w_n^{\bm{\xi}}) = \inf_{u\in U} H(t, x^*(t), u, p^*(t); w_n^{\bm{\xi}}) \quad \text{ for all } t \in [0,T].
\end{equation}
\end{theorem}

\begin{equation} 
\begin{aligned}
    H(t,x,u,p; w_n^{\bm{\xi}}) := &\langle p, b(x,u)\rangle + r(x, u) + \left( p^\top \sigma(x) - z_\theta(t,x) \right) \dot{w}_n^{\bm{\xi}}(t)  \\
    &+ \frac{1}{2}\op{Tr}(\sigma(x) J_xz_\theta(t,x)) - \langle p, \frac{1}{2}(\nabla \sigma:\sigma)(x)\rangle, \\
    \mathcal{H}(t,x,p;w_n^{\bm{\xi}}) :=&\inf_{u\in U} H(t,x,u,p; w_n^{\bm{\xi}}), 
\end{aligned}
\end{equation}

\begin{assumption} \label{assu:existence}
    Given any $\bm{\xi} \in \R^{nd}$. Assume that 
    \begin{enumerate}
        \item There exist positive constants $C_1, C_2$ such that 
        \begin{equation*}
            \begin{aligned}
                |\tilde{b}(t,x,u;w_n^{\bm{\xi}})| &\le C_1 (1 +|x| +|u|), \\
                |\tilde{b}(t,x',u;w_n^{\bm{\xi}}) - \tilde{b}(t,x,u;w_n^{\bm{\xi}})| &\le C_2|x-x'|(1 + |u|), 
            \end{aligned}
        \end{equation*}
        for all $t\in [0,T]$, $x, x'\in \R^d$, $u\in U$. 
        \item $U$ is closed. 
        \item For each $(t,x)\in [0,T]\times \R^d$, there exists a continuous function $\ell:U\to \R$ such that 
        \begin{equation*}
            \tilde{r}(t,x,u;w_n^{\bm{\xi}}) \ge \ell(u) \quad \text{ and }\quad \frac{\ell(u)}{|u|} \to +\infty \text{ as } |u|\to \infty. 
        \end{equation*}
        \item For each $(t,x)\in [0,T]\times \R^d$, the set 
        \begin{equation*}
            \tilde{F}(t,x;w_n^{\bm{\xi}}) := \{(z,z_{d+1})\in \R^{d+1}: z=\tilde{b}(t,x,u;w_n^{\bm{\xi}})\in \R^d, z_{d+1} \ge \tilde{r}(t,x,u;w_n^{\bm{\xi}})\in \R, u\in U\}
        \end{equation*}
        is convex. 
    \end{enumerate}
\end{assumption}

\begin{theorem}[Generalized Hopf formula] \label{thm:generalized-hopf}
Given any $\bm{\xi}^{(m)}\in \R^{nd}$. Let Assumptions~\ref{assu:dual-sc-prob}, ~\ref{assu:Pontryagin} and~\ref{assu:existence} hold. Assume that $g$ is proper convex and admits a convex conjugate $g^*$. Let $v_n^{(m)}:[0,T]\times \R^d \to \R$ be the solution to the problem~\eqref{eq:doc-prob} with $w_n = w_n^{\bm{\xi}^{(m)}}$. Then for any $x_0\in \R^d$, the value $v_n^{(m)}(0,x_0)$ can be represented as 
\begin{equation} \label{eq:generalized-hopf-formula}
\begin{aligned}
    v_n^{(m)}(0, x_0) = \sup_{p_0\in \R^d}\bigg\{ &\langle x_0, p_0\rangle - g^*(p(T)) + \int_0^T \bigg[ \mathcal{H}(t, x(t), p(t); w_n^{\bm{\xi}^{(m)}}) - \langle x(t), \partial_x \mathcal{H}(t, x(t), p(t); w_n^{\bm{\xi}^{(m)}})\rangle \bigg]\,\dd t: \\
    & \dot{x}(t) = \partial_p \mathcal{H}(t, x(t), p(t); w_n^{\bm{\xi}^{(m)}}), \quad x(0) = x_0, \\
    & \dot{p}(t) = -\partial_x \mathcal{H}(t,x(t),p(t); w_n^{\bm{\xi}^{(m)}}), \quad p(0)= p_0 \bigg\},  
\end{aligned}
\end{equation}
where $\mathcal{H}$ is defined in \eqref{eq:hamiltonian-defi}. 
\end{theorem}

\begin{proof}
{\bf Step 1.} We begin by characterizing the optimal cost using Pontryagin's Maximum Principle (Theorem \ref{thm:pontryagin}). The value function is given by 
\begin{equation} \label{eq:optimal-cost}
    v_n^{(m)}(0,x_0) = g(x^*(T)) + \int_0^T \tilde{r}(t,x^*(t), u^*(t);w_n^{\bm{\xi}^{(m)}})\,\dd t, 
\end{equation}
where the optimal trajectory $(x^*, p^*)$ solves the Hamiltonian system 
\begin{equation} \label{eq:Pontryagin-ode}
\begin{aligned}
    \dot{x}^*(t) &= \partial_p \mathcal{H}(t,x^*(t), p^*(t); w_n^{\bm{\xi}^{(m)}}), \quad &x^*(0) = x_0, \\
    \dot{p}^*(t) &= -\partial_x \mathcal{H}(t, x^*(t), p^*(t); w_n^{\bm{\xi}^{(m)}}), \quad &p^*(T) = \nabla g(x^*(T)). 
\end{aligned}
\end{equation}
Furthermore, by Assumption \ref{assu:existence} (existence of optimal control), for any $t\in [0,T]$ there exists $u^*(t)$ satisfying the minimization condition 
\begin{equation*}
    u^*(t) \in \operatorname*{arg\,inf}_{u\in U} H(t,x^*(t), u, p^*(t); w_n^{\bm{\xi}^{(m)}}).
\end{equation*}
We shall express the running cost $\tilde{r}$ in terms of the Hamiltonian and the state dynamics. 
Recalling the definition of $\mathcal{H}$ and observing from \eqref{eq:Pontryagin-ode} that $\dot{x}^*(t) = \tilde{b}(t,x^*(t), u^*(t); w_n^{\bm{\xi}^{(m)}})$, we have 
\begin{equation*}
    \mathcal{H}(t,x^*(t), p^*(t); w_n^{\bm{\xi}^{(m)}}) = \langle p^*(t), \tilde{b}(t,x^*(t), u^*(t); w_n^{\bm{\xi}^{(m)}})\rangle + \tilde{r}(t,x^*(t), u^*(t); w_n^{\bm{\xi}^{(m)}}). 
\end{equation*}
Rearranging terms yields the relationship 
\begin{equation} \label{eq:r-tilde-by-hamiltonian}
    \tilde{r}(t,x^*(t), u^*(t); w_n^{\bm{\xi}^{(m)}}) = \mathcal{H}(t,x^*(t), p^*(t); w_n^{\bm{\xi}^{(m)}}) - \langle p^*(t), \dot{x}^*(t) \rangle. 
\end{equation}

{\bf Step 2.} We establish the following identity for any $t\in [0,T]$, 
\begin{equation} \label{eq:equivalence-1}
    \mathcal{H}(t,x^*(t), p^*(t); w_n^{\bm{\xi}^{(m)}}) - \langle p^*(t), \dot{x}^*(t) \rangle = \sup_{p\in \R^d}\left\{ \mathcal{H}(t,x^*(t), p; w_n^{\bm{\xi}^{(m)}}) - \langle p, \dot{x}^*(t) \rangle \right\}.
\end{equation}
First, observe that for any fixed control $u\in U$, the mapping $p\mapsto H(t,x^*(t),u, p;w_n^{\bm{\xi}^{(m)}})$ is affine. Since $\mathcal{H}$ is defined as the infimum of these affine functions over $u$, the map $p\mapsto \mathcal{H}(t,x^*(t),p;w_n^{\bm{\xi}^{(m)}})$ is concave. 

Then, we define the convex function 
\begin{equation*}
    h(p) := -\mathcal{H}(t,x^*(t),p;w_n^{\bm{\xi}^{(m)}}). 
\end{equation*}
Its convex conjugate $h^*$ is given by 
\begin{equation*}
    h^*(q) := \sup_{p\in \mathbb{R}^d}\left\{ \langle p, q\rangle - h(p) \right\} = \sup_{p\in \R^d}\left\{ \langle p, q\rangle +\mathcal{H}(t,x^*(t), p;w_n^{\bm{\xi}^{(m)}}) \right\}
\end{equation*}
Evaluating $h^*$ at $q = -\dot{x}^*(t)$, we have 
\begin{equation}
\label{eq:hstar-xstar}
    h^*(-\dot{x}^*(t)) = \sup_{p\in \R^d} \left\{ \mathcal{H}(t,x^*(t), p;w_n^{\bm{\xi}^{(m)}}) - \langle p, \dot{x}^*(t)\rangle \right\}. 
\end{equation}
By the Fenchel-Young inequality, for any $p\in \R^d$, 
\begin{equation*}
    h^*(-\dot{x}^*(t)) + h(p) \ge \langle p, -\dot{x}^*(t) \rangle, 
\end{equation*}
with equality holding if and only if $-\dot{x}^*(t) =  -\partial_p \mathcal{H}(t,x^*(t),p;w_n^{\bm{\xi}^{(m)}})$. 
Recalling the state dynamics from Step 1, we have $\dot{x}^*(t) =  \partial_p \mathcal{H}(t,x^*(t),p;w_n^{\bm{\xi}^{(m)}})$. 
Thus, the equality condition is satisfied at $p = p^*(t)$ and we have 
\begin{equation*}
    h^*(-\dot{x}^*(t)) + h(p^*(t)) = \langle p^*(t), -\dot{x}^*(t) \rangle.
\end{equation*}
Rearranging this equality and substituting the definitions of $h$ and \eqref{eq:hstar-xstar} recovers equation~\eqref{eq:equivalence-1}. 

{\bf Step 3.} We now establish a lower bound for the optimal cost by passing to a dual formulation.
Since $g$ is proper convex and continuous by Assumption~\ref{assu:dual-sc-prob}, we express the terminal cost using its convex conjugate 
\begin{equation} \label{eq:conjugate-g}
    g(x^*(T)) = \sup_{p\in \R^d} \left\{ \langle x^*(T), p \rangle - g^*(p) \right\}. 
\end{equation}
Substituting \eqref{eq:conjugate-g} and the identity \eqref{eq:equivalence-1} from Step 2 into the cost equation \eqref{eq:optimal-cost}, we obtain 
\begin{equation*}
    v_n^{(m)}(0,x_0) = \sup_{p\in \R^d} \left\{ \langle x^*(T), p \rangle - g^*(p) \right\} + \int_0^T \sup_{p\in \R^d}\left\{ \mathcal{H}(t,x^*(t), p; w_n^{\bm{\xi}^{(m)}}) - \langle p, \dot{x}^*(t) \rangle \right\} \,\dd t 
\end{equation*}
We now relax the pointwise supremum inside the integral to a supremum over the space of measurable functions $p:[0,T]\to \R^d$: 
\begin{equation} \label{eq:sup-integral-inequality} 
v_n^{(m)}(0,x_0) \ge \sup_{\substack{p:[0,T]\to \R^d \\ \text{measurable}}} \left\{ \langle x^*(T), p(T) \rangle - g^*(p(T)) + \int_0^T \left[ \mathcal{H}(t,x^*(t), p(t); w_n^{\bm{\xi}^{(m)}}) - \langle p(t), \dot{x}^*(t) \rangle\right] \,\dd t \right\}. 
\end{equation}
Here, we make some remark on measurability. The interchange of supremum and integral is justified as follows. The map $(t,x,p,u)\mapsto \langle \tilde{b}(t,x,u;w_n^{\bm{\xi}^m}), p\rangle + \tilde{r}(t,x,p;w_n^{\bm{\xi}^m})$ is continuous by Assumption~\ref{assu:Pontryagin} and thus measurable. Since $U$ is a separable metric space, the infimum defining $\mathcal{H}$ can be taken over a countable dense subset, preserving measurability. Consequently, the map $t\mapsto \sup_{p\in \R^d}\{\mathcal{H}(t, x^*(t), p;w_n^{\bm{\xi}^{(m)}}) - \langle p, \dot{x}^*(t)\rangle\}$ is measurable. Its integrability is guaranteed by the integrability of $\tilde{r}$ and the equality derived in \eqref{eq:r-tilde-by-hamiltonian} and \eqref{eq:equivalence-1}. Moreover, the measurability of $t\mapsto \mathcal{H}(t,x^*(t), p(t); w_n^{\bm{\xi}^{(m)}}) - \langle p(t), \dot{x}^*(t) \rangle$ is ensured since we restrict $p:[0,T]\to \R^d$ to be measurable. 

Next, applying integration by parts to the term involving $\dot{x}^*(t)$ yields 
\begin{equation} \label{eq:optimize-over-all-p}
\begin{aligned}
    v_n^{(m)}(0,x_0) &\ge \sup_{\substack{p:[0,T]\to \R^d \\ \text{measurable}}} \bigg\{ x^*(T), p(T) \rangle - g^*(p(T)) + \int_0^T \mathcal{H}(t,x^*(t), p(t); w_n^{\bm{\xi}^{(m)}})\,\dd t \\
    &\qquad \qquad \qquad - \langle x^*(T), p(T)\rangle + \langle x_0, p(0) \rangle + \int_0^T \langle x^*(t), \dot{p}(t) \rangle \,\dd t \bigg\} \\
    &= \sup_{\substack{p:[0,T]\to \R^d \\ \text{measurable}}} \bigg\{\langle x_0, p(0) \rangle - g^*(p(T)) + \int_0^T \left(\mathcal{H}(t,x^*(t), p(t); w_n^{\bm{\xi}^{(m)}}) + \langle x^*(t), \dot{p}(t) \rangle\right) \,\dd t \bigg\}. 
\end{aligned}
\end{equation}
Then, we intend to optimize over a subset of $\{p:[0,T]\to \R^d\mid p \text{ is measurable}\}$. Define the set $S_p^{x_0}$ of continuous functions $p(\cdot)$ that are solutions to the Hamiltonian system by 
\begin{equation*}
S_p^{x_0} := \left\{
p\in C([0,T]; \mathbb{R}^d) \;\bigg|\; 
\exists x: [0,T] \to \mathbb{R}^d \text{ s.t. }
\begin{array}{l}
\dot{p}(t) = -\partial_x \mathcal{H}(t,x(t), p(t); w_n^{\bm{\xi}^{(m)}}), \\
\dot{x}(t) = \partial_p\mathcal{H}(t,x(t), p(t); w_n^{\bm{\xi}^{(m)}}), \quad x(0) = x_0
\end{array}
\right\}.
\end{equation*}
Here, $p(0)$ is ``free". By Assumption~\ref{assu:existence} (existence of optimal control) and Theorem~\ref{thm:pontryagin}, there exists an optimal trajectory $(x^*, p^*)$ such that $p^* \in S_p^{x_0}$ and thus $S_p^{x_0}$ is nonempty.
Restricting the domain of the supremum in \eqref{eq:optimize-over-all-p} to $S_p^{x_0}$ yields the lower bound: 
\begin{equation} \label{eq:quantity-I}
    v_n^{(m)}(0,x_0) \ge \sup_{p\in S_p^{x_0}} \bigg\{\langle x_0, p(0) \rangle - g^*(p(T)) + \int_0^T \left(\mathcal{H}(t,x^*(t), p(t); w_n^{\bm{\xi}^{(m)}}) + \langle x^*(t), \dot{p}(t) \rangle\right) \,\dd t \bigg\} =: \textrm{I}. 
\end{equation}

{\bf Step 4.} Next, we will show that 
\begin{equation} \label{eq:identity-I-II}
    \textrm{I} \ge \textrm{II} := \sup_{(x,p)\in S_{xp}^{x_0}} \bigg\{\langle x_0, p(0) \rangle - g^*(p(T)) + \int_0^T \left(\mathcal{H}(t,x(t), p(t); w_n^{\bm{\xi}^{(m)}}) + \langle x(t), \dot{p}(t) \rangle\right) \,\dd t \bigg\}, 
\end{equation}
where the admissible set $S_{xp}^{x_0}$ is defined by 
\begin{equation*}
    S_{xp}^{x_0} := \left\{
(x,p)\in C([0,T]; \mathbb{R}^d)^2 \;\bigg|\;
\begin{array}{l}\dot{p}(t) = -\partial_x \mathcal{H}(t,x(t), p(t); w_n^{\bm{\xi}^{(m)}}), \\
\dot{x}(t) = \partial_p\mathcal{H}(t,x(t), p(t); w_n^{\bm{\xi}^{(m)}}),\quad x(0) = x_0
\end{array}\right\}. 
\end{equation*}
For any $(x, p) \in S_{xp}^{x_0}$, it holds that $p\in S_p^{x_0}$. Therefore, to establish \eqref{eq:identity-I-II}, it suffices to prove that the integrand satisfy the pointwise inequality 
\begin{equation} \label{eq:goal-ineq}
    \mathcal{H}(t, x^*(t), p(t); w_n^{\bm{\xi}^{(m)}}) + \langle x^*(t), \dot{p}(t)\rangle \ge \mathcal{H}(t,x(t), p(t); w_n^{\bm{\xi}^{(m)}}) + \langle x(t), \dot{p}(t) \rangle. 
\end{equation}
Define the convex conjugate of $\mathcal{H}$ with respect to $x$ by 
\begin{equation*}
    F(\eta; t, p(t), w_n^{\bm{\xi}^{(m)}}) := \sup_{x\in \R^d} \left\{ \langle \eta, x\rangle - \mathcal{H}(t, x, p(t); w_n^{\bm{\xi}^{(m)}}) \right\}. 
\end{equation*}
By the Fenchel–Young inequality, evaluating the conjugate at $\eta = -\dot{p}(t)$ and the primal function at $x^*(t)$ yields 
\begin{equation} \label{eq:FY-1}
    \mathcal{H}(t, x^*(t), p(t); w_n^{\bm{\xi}^{(m)}}) + F(-\dot{p}(t); t, p(t), w_n^{\bm{\xi}^{(m)}}) \ge \langle x^*(t), -\dot{p}(t)\rangle. 
\end{equation}
Furthermore, the Fenchel-Young inequality becomes an equality if $\eta = \partial_x\mathcal{H}$. Since $(x, p) \in S_{xp}^{x_0}$ implies that $-\dot{p}(t) = \partial_x\mathcal{H}(t, x(t), p(t); w_n^{\bm{\xi}^{(m)}})$, we obtain the identity 
\begin{equation} \label{eq:FY-2}
    \mathcal{H}(t, x(t), p(t); w_n^{\bm{\xi}^{(m)}}) + F(-\dot{p}(t); t, p(t), w_n^{\bm{\xi}^{(m)}}) = \langle x(t), -\dot{p}(t)\rangle. 
\end{equation}
Combining \eqref{eq:FY-1} and \eqref{eq:FY-2} leads to the pointwise inequality \eqref{eq:goal-ineq}, which in turn confirms $\textrm{I} \ge \textrm{II}$.
Therefore, by \eqref{eq:quantity-I} and \eqref{eq:identity-I-II}, we obtain 
\begin{equation} \label{eq:generalized-hopf-one-side}
\begin{aligned}
    v_n^{(m)}(0, x_0) \ge \sup_{p_0\in \R^d}\bigg\{ &\langle x_0, p_0\rangle - g^*(p(T)) + \int_0^T \bigg[ \mathcal{H}(t, x(t), p(t); w_n^{\bm{\xi}^{(m)}}) + \langle x(t), \dot{p}(t)\rangle \bigg]\,\dd t: \\
    & \dot{x}(t) = \partial_p \mathcal{H}(t, x(t), p(t); w_n^{\bm{\xi}^{(m)}}), \quad x(0) = x_0, \\
    & \dot{p}(t) = -\partial_x \mathcal{H}(t,x(t),p(t); w_n^{\bm{\xi}^{(m)}}), \quad p(0)= p_0 \bigg\}. 
\end{aligned}
\end{equation}

{\bf Step 5.} Finally, it remains to show that the inequality in \eqref{eq:generalized-hopf-one-side} is in fact an equality. We use the Fenchel-Young inequality 
\begin{equation*}
    \langle x(T), p(T)\rangle \le g(x(T)) + g^*(p(T)), 
\end{equation*}
where the equality holds only when $p(T) = \nabla g(x(T))$, since $g$ is continuously differentiable by Assumption~\ref{assu:Pontryagin}. By \eqref{eq:Pontryagin-ode}, we have $(x^*, p^*)\in S_{xp}^{x_0}$ and $p^*(T) = \nabla g(x^*(T))$.
Thus, there holds 
\begin{equation} \label{eq:equality-g-g-star}
    \langle x^*(T), p^*(T)\rangle = g(x^*(T)) + g^*(p^*(T)). 
\end{equation}
Using \eqref{eq:generalized-hopf-one-side}, \eqref{eq:equality-g-g-star}, and integration by parts, we obtain 
\begin{equation*}
    \begin{aligned}
        v_n^{(m)}(0,x_0) &\ge \langle x_0, p^*(0)\rangle - g^*(p^*(T)) + \int_0^T \bigg[ \mathcal{H}(t, x^*(t), p^*(t); w_n^{\bm{\xi}^{(m)}}) + \langle x^*(t), \dot{p}^*(t)\rangle \bigg]\,\dd t \\
        &= \langle x_0, p^*(0)\rangle - g(x^*(T)) - \langle x^*(T), p^*(T)\rangle + \int_0^T \bigg[\mathcal{H}(t, x^*(t), p^*(t); w_n^{\bm{\xi}^{(m)}})\,\dd t + \langle x^*(t), \dot{p}^*(t)\rangle \bigg]\,\dd t \\
        &= g(x^*(T)) + \int_0^T \bigg[\mathcal{H}(t, x^*(t), p^*(t); w_n^{\bm{\xi}^{(m)}})\,\dd t + \langle p^*(t), \dot{x}^*(t)\rangle \bigg]\,\dd t \\
        &= v_n^{(m)}(0,x_0),
    \end{aligned}
\end{equation*}
where the last equality holds due to \eqref{eq:optimal-cost} and \eqref{eq:r-tilde-by-hamiltonian}. 
This verifies that \eqref{eq:generalized-hopf-one-side} is an equality and thus completes the proof. 
\end{proof}

\end{document}
